% Part: proof-theory
% Chapter: sequent-calculus
% Section: rules-G3c

\documentclass[../../../include/open-logic-section]{subfiles}

\begin{document}

\begin{table}
\begin{tabular}{@{}c@{}c@{}}
\multicolumn{2}{@{}c@{}}{公理}\\[2ex]
$!C, \Gamma \Sequent \Delta, !C$ & $\lfalse, \Gamma \Sequent \Delta$
\\[2ex]
\multicolumn{2}{@{}c@{}}{逻辑规则}\\[2ex]
\Axiom$\lnot !A, \Gamma \fCenter \Delta, !A $
\RightLabel{\LeftR{\lnot}}
\UnaryInf$\lnot !A, \Gamma \fCenter \Delta$
\DisplayProof
&
\Axiom$!A, \Gamma \fCenter$
\RightLabel{\RightR{\lnot}}
\UnaryInf$ \Gamma \fCenter \Delta, \lnot !A $
\DisplayProof
\\[4ex]
\Axiom$ !A, !B, \Gamma \fCenter \Delta$
\RightLabel{\LeftR{\land}}
\UnaryInf$ !A \land !B, \Gamma \fCenter \Delta$
\DisplayProof
&
\Axiom$\Gamma \fCenter \Delta, !A$
\Axiom$ \Gamma \fCenter \Delta, !B$
\RightLabel{\RightR{\land}}
\BinaryInf$ \Gamma \fCenter \Delta, !A \land !B $
\DisplayProof
\\[4ex]\Axiom$!A, \Gamma \fCenter \Delta$
\Axiom$!B, \Gamma \fCenter \Delta$
\RightLabel{\LeftR{\lor}}
\BinaryInf$!A \lor !B, \Gamma \fCenter \Delta$
\DisplayProof
&
\Axiom$\Gamma \fCenter \Delta, !A, !B$
\RightLabel{\RightR{\lor}}
\UnaryInf$ \Gamma \fCenter \Delta, !A \lor !B$
\DisplayProof
\\[4ex]
\Axiom$!A \lif !B, \Gamma \fCenter \Delta, !A$
\Axiom$ !B, \Gamma \fCenter \Delta$
\RightLabel{\LeftR{\lif}}
\BinaryInf$!A \lif !B, \Gamma \fCenter \Delta$
\DisplayProof
&
\Axiom$ !A, \Gamma \fCenter !B$
\RightLabel{\RightR{\lif}}
\UnaryInf$ \Gamma \fCenter \Delta, !A \lif !B $
\DisplayProof
\\[4ex]
\Axiom$ !A(t), \lforall[x][!A(x)]\Gamma \fCenter \Delta$
\RightLabel{\LeftR{\lforall}}
\UnaryInf$ \lforall[x][!A(x)],\Gamma \fCenter \Delta$
\DisplayProof
&
\Axiom$ \Gamma \fCenter !A(a) $
\RightLabel{\RightR{\lforall}}
\UnaryInf$ \Gamma \fCenter \lforall[x][!A(x)]$
\DisplayProof
\\[4ex]
\Axiom$ !A(a), \Gamma \fCenter \Delta $
\RightLabel{\LeftR{\lexists}}
\UnaryInf$ \lexists[x][!A(x)], \Gamma \fCenter \Delta$
\DisplayProof
&
\Axiom$ \Gamma \fCenter \Delta, \lexists[x][!A(x)], !A(t) $
\RightLabel{\RightR{\lexists}}
\UnaryInf$ \Gamma \fCenter \Delta, \lexists[x][!A(x)]$
\DisplayProof
\end{tabular}
\caption{直觉主义逻辑的多结论相继式演算 \Log{mG3i} 的规则。在 $\RightR{\forall}$ 和 $\LeftR{\exists}$ 中，常项符号~$a$ 不得出现在结论相继式中。$\Gamma$ 和 $\Delta$ 是公式的多重集。在公理中，公式~$!C$ 必须是原子的。\Log{mG1m} 是去掉公理 $\lfalse,
\Gamma \Sequent \Delta$ 后的 \Log{mG1i} 系统。}
\ollabel{tab:G3c}
\end{table}

\end{document}